\bookchapter{Kapitel 53: Datenspeicherung und Netzwerkanfragen}

\chapterintro{Eine App, die nur “angucken” kann, ist tot; erst eine App, die Daten speichern und sich vernetzen kann, lebt. In diesem Kapitel lernst du das “Trio” der Android-Datenschicht kennen: DataStore (kleine Einstellungen speichern), Room (strukturierte Daten speichern), Retrofit (mit dem Server sprechen) – plus ihr “Dispositionszentrum”, die Koroutinen. Nach diesem Kapitel kannst du echte Apps schreiben, die sich Nutzer-Gewohnheiten merken und Inhalte live aktualisieren.}

\booksection{Welches Speicherkonzept wofür? Erst diese Tabelle anschauen}

\begin{center}
\begin{tabularx}{\linewidth}{|X|X|X|X|}
\hline
\textbf{Konzept} & \textbf{Wofür geeignet?} & \textbf{Eigenschaften} & \textbf{Wo im Buch genutzt?} \tabularnewline
\hline
DataStore & Verstreute Einstellungen (Theme, Schalter) & Schlüssel-Wert-Paare, einfach und sicher & Kapitel 54: Einstellungen des Haushaltsbuchs \tabularnewline
Room & Strukturierte Daten (Buchungen, Nutzer) & Eine echte SQLite-Datenbank verpackt, abfragbar & Kapitel 54: Buchungseinträge \tabularnewline
Retrofit & Serverdaten (nicht speichern, online abrufen) & Der De-facto-Standard unter den HTTP-Bibliotheken & Kapitel 55: Wetterdaten \tabularnewline
\hline
\end{tabularx}
\end{center}

\begin{center}\small\color{gray}Die Arbeitsteilung des Daten-Trios\end{center}

Merksatz: “Kleine Einstellungen ins DataStore, Daten in Mengen ins Room, Daten von draußen über Retrofit.” Das in alten Tutorials allgegenwärtige SharedPreferences ist von DataStore abgelöst, handgeschriebenes SQLite-SQL von Room – lern gleich das Neue, geh nicht den alten Weg.

\booksection{DataStore: Kleine Einstellungen wegspeichern}

Werte wie “die zuletzt gewählte Währung” oder “ist der Dunkelmodus an?” willst du verstreut abspeichern – dafür ist DataStore (Preferences-Variante) da:

\begin{lstlisting}[language=Kotlin]
val Context.dataStore by preferencesDataStore(name = "settings")

// schreiben: Wert speichern
suspend fun saveTheme(isDark: Boolean) {
    context.dataStore.edit { prefs ->
        prefs[THEME_KEY] = isDark
    }
}

// lesen: herausholen (Flow meldet Änderungen automatisch)
val themeFlow: Flow<Boolean> = context.dataStore.data
    .map { it[THEME_KEY] ?: false }
\end{lstlisting}

Achte darauf: Zurückgegeben wird ein Flow – ein Datenstrom, der “fließt”. Sobald sich der Wert ändert, bekommen alle Stellen, die ihn beobachten, automatisch den neuen Wert. Dieses Konzept (Reaktivität) ist der Grundstein des modernen Android – in Kapitel 54 kommt es zum Praxiseinsatz.

\booksection{Room: Den Daten ein “richtiges Zuhause” geben}

Ein Haushaltsbuch hat hunderte, tausende Buchungen, die nach Kategorie und Monat abgefragt und ausgewertet werden sollen – dafür braucht es eine Datenbank. Room ist die offizielle Hülle um SQLite und macht “Tabellen anlegen, Daten abfragen” so natürlich wie Kotlin-Schreiben. Drei Dinge:

\begin{lstlisting}[language=Kotlin]
// 1. Entity: die Tabellenstruktur (eine Tabelle = eine Data Class)
@Entity(tableName = "records")
data class Record(
    @PrimaryKey(autoGenerate = true) val id: Long = 0,
    val amount: Double,
    val category: String,
    val note: String,
    val date: String
)

// 2. Dao: das Operations-Interface (anlegen, löschen, ändern, lesen)
@Dao
interface RecordDao {
    @Query("SELECT * FROM records ORDER BY date DESC")
    fun getAll(): Flow<List<Record>>

    @Insert
    suspend fun insert(record: Record)

    @Delete
    suspend fun delete(record: Record)

    @Query("SELECT SUM(amount) FROM records WHERE category = :c")
    suspend fun sumByCategory(c: String): Double
}

// 3. Database: die Endmontage
@Database(entities = [Record::class], version = 1)
abstract class AppDatabase : RoomDatabase() {
    abstract fun recordDao(): RecordDao
}
\end{lstlisting}

Dieses “Trio”-Template kannst du unverändert in jedes deiner Projekte übernehmen – einfach die Feldnamen austauschen. Die Abfragen liefern Flow, die Oberfläche aktualisiert sich automatisch mit den Daten – herrlich.

\booksection{Retrofit: Mit dem Server sprechen}

Wetter, Nachrichten, Wechselkurse … Diese Daten liegen nicht auf dem Telefon, sondern auf den Servern von anderen. Retrofit macht “Daten vom Server anfordern” so einfach wie einen Funktionsaufruf. In drei Schritten:

\begin{lstlisting}[language=Kotlin]
// 1. Interface definieren: HTTP-Anfragen als Funktionen beschreiben
interface WeatherApi {
    @GET("v1/weather")
    suspend fun getWeather(@Query("city") city: String): WeatherResponse
}

// 2. Instanz erzeugen
val api = Retrofit.Builder()
    .baseUrl("https://api.example.com/")
    .addConverterFactory(GsonConverterFactory.create())
    .build()
    .create(WeatherApi::class.java)

// 3. aufrufen wie eine Funktion (in einer Koroutine)
val result = api.getWeather("Berlin")
\end{lstlisting}

\textbf{Stolperfalle: Zwei eiserne Netzwerk-Regeln}\par

(1) Netzwerkanfragen dürfen niemals im Hauptthread laufen – die Oberfläche friert ein oder die App stürzt ab (NetworkOnMainThreadException). Alle Netzwerkaufrufe gehören in eine Koroutine (viewModelScope.launch). (2) Klartext-HTTP ist standardmäßig verboten: Ab Android 9 wird HTTPS verlangt. Für Tests mit lokalen http://-Adressen musst du im Manifest usesCleartextTraffic explizit deklarieren – ein echtes Produkt muss aber komplett über HTTPS laufen.

\booksection{JSON-Parsen: Den Netzdaten eine “Übersetzung” geben}

Was der Server zurückgibt, ist meist JSON – ein Text wie ein “verschachteltes Wörterbuch”. Retrofit braucht einen “Dolmetscher”, der daraus Kotlin-Objekte macht. 2026 gibt es zwei Mainstream-Lösungen:

\begin{center}
\begin{tabularx}{\linewidth}{|X|X|X|}
\hline
\textbf{Konzept} & \textbf{Eigenschaften} & \textbf{Empfohlen für} \tabularnewline
\hline
Gson / Moshi & Der klassische Reflection-Ansatz, null Konfigurationsaufwand & Schneller Start, kleine Projekte mit wenigen Feldern \tabularnewline
kotlinx.serialization & Die offizielle Lösung zur Compiler-Zeit, schneller und sicherer & Erste Wahl für neue Projekte (2026er Mainstream) \tabularnewline
\hline
\end{tabularx}
\end{center}

\begin{center}\small\color{gray}Die beiden großen JSON-Lösungen\end{center}

\begin{lstlisting}[language=Kotlin]
// versieh die Data Class mit @Serializable
@Serializable
data class WeatherResponse(
    val weather: List<Weather>,
    val main: Main,
    val name: String
)

// Retrofit bekommt den passenden Converter angehängt (Beispiel mit offizieller Unterstützung ab Retrofit 2.11+):
// .addConverterFactory(json.asConverterFactory("application/json".toMediaType()))
\end{lstlisting}

Gleich welche Lösung – das eiserne Gesetz lautet: Erst das echte zurückgegebene JSON anschauen, dann die Felder definieren. Viele “Parsen fehlgeschlagen”-Bugs entstehen, weil Feldnamen verschrieben oder Verschachtelungsebenen falsch geraten sind – gib das rohe JSON aus und gleiche ab, in zehn Sekunden ist der Fehler geortet.

\booksection{Koroutinen: Die “Gemeinsprache” der Android-Concurrency}

\booksection{Koroutinen: Die “Gemeinsprache” der Android-Concurrency}

Das immer wieder auftauchende suspend und launch stammen aus den Kotlin-Koroutinen – der offiziellen Concurrency-Lösung von Android. Grob gesagt: Operationen, die “ein Weilchen dauern” (Netzwerk, Datenbank), blockieren die Oberfläche nicht.

\begin{lstlisting}[language=Kotlin]
// starte eine Koroutine im ViewModel
viewModelScope.launch {
    val records = dao.getAll().first()  // suspend: hält automatisch an und wartet
    // Zustand der Oberfläche aktualisieren …
}

// Spickzettel der Schlüsselwörter:
// suspend   —— "ich muss warten, halte andere nicht auf"
// launch    —— "auf geht's!"
// viewModelScope —— die "Werkstatt", die der Ansicht folgt; räumt sich beim Zerstören selbst ab
\end{lstlisting}

Als Einsteiger musst du zunächst nur die Anwendung kennen; das Prinzip dahinter (“Anhalten und Fortsetzen”) kannst du später vertiefen. 90 \% dessen, was du im Alltag schreibst, ist diese eine Schablone.

\begin{tcbpractice}{Selbst ausprobieren: Eine kleine App "Satz des Tages" bauen}
Such dir eine kostenlose Zitate-API (z. B. hitokoto o. Ä.), hole mit Retrofit den “Satz des Tages” und zeige ihn mitten auf dem Bildschirm; füge einen Button “Noch einen” hinzu (neue Anfrage in einer Koroutine). Speichere Lieblingssätze mit Room in eine “Favoriten”-Sammlung. Diese Übung nutzt DataStore/Room/Retrofit/Korouten je einmal – danach sind Kapitel 7 und 8 ein Kinderspiel.

\end{tcbpractice}

\begin{tcbsummary}{Zusammenfassung des Kapitels}
\begin{itemize}
  \item Merksatz: kleine Einstellungen ins DataStore, Daten in Mengen ins Room, Daten von draußen über Retrofit
  \item DataStore liefert Flow: Ändert sich der Wert, bekommen Beobachter ihn automatisch
  \item Room-Trio: @Entity (Tabelle) + @Dao (Operationen) + @Database (Endmontage) – die Schablone ist direkt wiederverwendbar
  \item Retrofit in drei Schritten: Interface definieren Instanz erzeugen in der Koroutine aufrufen
  \item Zwei eiserne Regeln: Netzwerk nie im Hauptthread; echte Produkte immer über HTTPS
  \item Koroutine = Warten, ohne die Oberfläche zu blockieren; viewModelScope ist die eigene Werkstatt jeder Seite
\end{itemize}

\end{tcbsummary}
